% GENERATED FILE. Edit canonical lessons, then run npm run generate:books.
% canonical-source-hash: fnv1a64:c80ec734b9a16043
% canonical-lessons: JA-C309-kureru, JA-C309-ageru, JA-C309-okaeshi, JA-C309-hito, JA-C309-uchiakeru

\chapter{Presents in the Post}
\label{ch:ja-presents-in-the-post}
\hlchaptermodality{309}

\begin{chapteropening}
\textbf{By the end of this chapter:} \emph{I can say who gave what to whom.}

Complete the last lesson of chapter 309: I can say who gave what to whom.
\end{chapteropening}

\section[\texorpdfstring{\ja{くれる} (kureru)}{くれる (kureru)}]{\ja{くれる} (kureru) — to give to me; to give (to my family)}

\label{lesson:JA-C309-kureru}



\begin{quote}
Before the new one: say the Japanese for to speak to (someone), then the Japanese for busy.
\end{quote}

\subsection*{You'll want to know: \ja{くれる}}
\textbf{\ja{くれる}} — \emph{kureru} — \textquotedblleft{}to give to me; to give (to my family)\textquotedblright{}.

\begin{quote}
\textbf{\ja{そぼが} \ja{ゆかたを} \ja{くれました。}} — \emph{sobo ga yukata o kuremashita} — \textquotedblleft{}My grandmother gave me a summer kimono.\textquotedblright{}
\end{quote}

Polite \textbf{\ja{くれます}}, \emph{kuremasu}; an ichidan verb. Past (polite): \textbf{\ja{ともだちが} \ja{てがみを} \ja{くれました。}}, \emph{tomodachi ga tegami o kuremashita}.

\begin{cousinweb}
\textbf{\ja{くれる}} (\emph{kureru}) is a native Japanese word (wago), inherited rather than borrowed from Chinese, and it is usually written in kana (JMdict).
\end{cousinweb}

\begin{grammarlens}[title={What you've built}]
One more word for telling what happened.
\end{grammarlens}

\subsection*{Your turn}
\begin{itemize}
\raggedright
  \item \emph{Say it:} \emph{kureru}
  \item \emph{Say it:} \emph{kureru}, once more
  \item \emph{Say it:} \emph{sobo ga yukata o kuremashita} — My grandmother gave me a summer kimono.
  \item \emph{Recall:} say the Japanese for to speak to (someone), then the Japanese for busy, then say \emph{kureru} again
\end{itemize}

\begin{tcolorbox}[breakable,colback=teal!4,colframe=teal!35!black,title={Before you move on}]
What does \ja{くれる} mean? (\textquotedblleft{}To give to me; to give (to my family)\textquotedblright{}.) Say it once more.
\end{tcolorbox}
\section[\texorpdfstring{\ja{あげる} (ageru)}{あげる (ageru)}]{\ja{あげる} (ageru) — to give (to someone else)}

\label{lesson:JA-C309-ageru}



\begin{quote}
Before the new one: say the Japanese for busy, then the Japanese for to give to me.
\end{quote}

\subsection*{You'll want to know: \ja{あげる}}
\textbf{\ja{あげる}} — \emph{ageru} — \textquotedblleft{}to give (to someone else)\textquotedblright{}.

\begin{quote}
\textbf{\ja{いもうとに} \ja{かばんを} \ja{あげました。}} — \emph{imōto ni kaban o agemashita} — \textquotedblleft{}I gave my little sister a bag.\textquotedblright{}
\end{quote}

Polite \textbf{\ja{あげます}}, \emph{agemasu}; an ichidan verb. Past (polite): \textbf{\ja{あねが} \ja{こどもに} \ja{おかしを} \ja{あげました。}}, \emph{ane ga kodomo ni okashi o agemashita}.

\begin{cousinweb}
\textbf{\ja{あげる}} (\emph{ageru}) is a native Japanese word (wago), inherited rather than borrowed from Chinese; its kanji was attached later as a kun'yomi (JMdict).
\end{cousinweb}

\begin{grammarlens}[title={What you've built}]
One more word for telling what happened.
\end{grammarlens}

\subsection*{Your turn}
\begin{itemize}
\raggedright
  \item \emph{Say it:} \emph{ageru}
  \item \emph{Say it:} \emph{ageru}, once more
  \item \emph{Say it:} \emph{imōto ni kaban o agemashita} — I gave my little sister a bag.
  \item \emph{Recall:} say the Japanese for busy, then the Japanese for to give to me, then say \emph{ageru} again
\end{itemize}

\begin{tcolorbox}[breakable,colback=teal!4,colframe=teal!35!black,title={Before you move on}]
What does \ja{あげる} mean? (\textquotedblleft{}To give (to someone else)\textquotedblright{}.) Say it once more.
\end{tcolorbox}
\section[\texorpdfstring{\ja{おかえし} (okaeshi)}{おかえし (okaeshi)}]{\ja{おかえし} (okaeshi) — a gift in return}

\label{lesson:JA-C309-okaeshi}



\begin{quote}
Before the new one: say the Japanese for to give to me, then the Japanese for to give (to someone else).
\end{quote}

\subsection*{You'll want to know: \ja{おかえし}}
\textbf{\ja{おかえし}} — \emph{okaeshi} — \textquotedblleft{}a gift in return\textquotedblright{}.

\begin{quote}
\textbf{\ja{おいわいの} \ja{おかえしを} \ja{おくりました。}} — \emph{oiwai no okaeshi o okurimashita} — \textquotedblleft{}I sent a gift in return for the present.\textquotedblright{}
\end{quote}

\begin{cousinweb}
\textbf{\ja{おかえし}} (\emph{okaeshi}) is a native Japanese word (wago), inherited rather than borrowed from Chinese, built from \emph{o-}, \textquotedblleft{}the polite prefix\textquotedblright{}, and \emph{kaeshi}, \textquotedblleft{}returning\textquotedblright{} (from \emph{kaesu}) (JMdict, KANJIDIC2).
\end{cousinweb}

\begin{grammarlens}[title={What you've built}]
One more word for telling what happened.
\end{grammarlens}

\subsection*{Your turn}
\begin{itemize}
\raggedright
  \item \emph{Say it:} \emph{okaeshi}
  \item \emph{Say it:} \emph{okaeshi}, once more
  \item \emph{Say it:} \emph{oiwai no okaeshi o okurimashita} — I sent a gift in return for the present.
  \item \emph{Recall:} say the Japanese for to give to me, then the Japanese for to give (to someone else), then say \emph{okaeshi} again
\end{itemize}

\begin{tcolorbox}[breakable,colback=teal!4,colframe=teal!35!black,title={Before you move on}]
What does \ja{おかえし} mean? (\textquotedblleft{}A gift in return\textquotedblright{}.) Say it once more.
\end{tcolorbox}
\section[\texorpdfstring{\ja{ひと} (hito)}{ひと (hito)}]{\ja{ひと} (hito) — a person; people}

\label{lesson:JA-C309-hito}



\begin{quote}
Before the new one: say the Japanese for to give (to someone else), then the Japanese for a gift in return.
\end{quote}

\subsection*{You'll want to know: \ja{ひと}}
\textbf{\ja{ひと}} — \emph{hito} — \textquotedblleft{}a person; people\textquotedblright{}.

\begin{quote}
\textbf{\ja{となりの} \ja{ひとが} \ja{やさいを} \ja{くれました。}} — \emph{tonari no hito ga yasai o kuremashita} — \textquotedblleft{}The person next door gave us some vegetables.\textquotedblright{}
\end{quote}

\begin{cousinweb}
\textbf{\ja{ひと}} (\emph{hito}) is a native Japanese word (wago), inherited rather than borrowed from Chinese; its kanji was attached later as a kun'yomi (JMdict). It is the \emph{-bito} of \emph{koibito}, \textquotedblleft{}a sweetheart\textquotedblright{}.
\end{cousinweb}

\begin{grammarlens}[title={What you've built}]
One more word for telling what happened.
\end{grammarlens}

\subsection*{Your turn}
\begin{itemize}
\raggedright
  \item \emph{Say it:} \emph{hito}
  \item \emph{Say it:} \emph{hito}, once more
  \item \emph{Say it:} \emph{tonari no hito ga yasai o kuremashita} — The person next door gave us some vegetables.
  \item \emph{Recall:} say the Japanese for to give (to someone else), then the Japanese for a gift in return, then say \emph{hito} again
\end{itemize}

\begin{tcolorbox}[breakable,colback=teal!4,colframe=teal!35!black,title={Before you move on}]
What does \ja{ひと} mean? (\textquotedblleft{}A person; people\textquotedblright{}.) Say it once more.
\end{tcolorbox}
\section[\texorpdfstring{\ja{うちあける} (uchiakeru)}{うちあける (uchiakeru)}]{\ja{うちあける} (uchiakeru) — to confide; to open up}

\label{lesson:JA-C309-uchiakeru}



\begin{quote}
Before the new one: say the Japanese for a gift in return, then the Japanese for a person.
\end{quote}

\subsection*{You'll want to know: \ja{うちあける}}
\textbf{\ja{うちあける}} — \emph{uchiakeru} — \textquotedblleft{}to confide; to open up\textquotedblright{}.

\begin{quote}
\textbf{\ja{ともだちに} \ja{ひみつを} \ja{うちあけました。}} — \emph{tomodachi ni himitsu o uchiakemashita} — \textquotedblleft{}I told my friend my secret.\textquotedblright{}
\end{quote}

Polite \textbf{\ja{うちあけます}}, \emph{uchiakemasu}; an ichidan verb. Past (polite): \textbf{\ja{あねが} \ja{ひみつを} \ja{うちあけました。}}, \emph{ane ga himitsu o uchiakemashita}.

\begin{cousinweb}
\textbf{\ja{うちあける}} (\emph{uchiakeru}) is a native Japanese word (wago), inherited rather than borrowed from Chinese, built from \emph{uchi}, \textquotedblleft{}striking\textquotedblright{} (from \emph{utsu}), and \emph{akeru}, \textquotedblleft{}to open\textquotedblright{} (JMdict, KANJIDIC2).
\end{cousinweb}

\begin{grammarlens}[title={What you've built}]
One more word for telling what happened.
\end{grammarlens}

\subsection*{Your turn}
\begin{itemize}
\raggedright
  \item \emph{Say it:} \emph{uchiakeru}
  \item \emph{Say it:} \emph{uchiakeru}, once more
  \item \emph{Say it:} \emph{tomodachi ni himitsu o uchiakemashita} — I told my friend my secret.
  \item \emph{Recall:} say the Japanese for a gift in return, then the Japanese for a person, then say \emph{uchiakeru} again
\end{itemize}

\begin{tcolorbox}[breakable,colback=teal!4,colframe=teal!35!black,title={Before you move on}]
What does \ja{うちあける} mean? (\textquotedblleft{}To confide; to open up\textquotedblright{}.) Say it once more.
\end{tcolorbox}
